\documentclass[11pt]{article}

\usepackage[margin=1in]{geometry}
\usepackage[T1]{fontenc}
\usepackage{lmodern}
\usepackage{amsmath,amssymb,amsthm,mathtools}
\usepackage{enumitem}
\usepackage[colorlinks=true,linkcolor=blue,urlcolor=blue]{hyperref}

\setlist[itemize]{topsep=2pt,itemsep=2pt}
\setlist[enumerate]{topsep=2pt,itemsep=2pt}
\setlength{\parskip}{0.5em}
\setlength{\parindent}{0pt}

\newcommand{\F}{\mathbb F}
\newcommand{\R}{\mathbb R}
\newcommand{\C}{\mathbb C}
\newcommand{\Mat}{\operatorname{Mat}}
\newcommand{\id}{\operatorname{id}}
\newcommand{\tr}{\operatorname{tr}}
\newcommand{\rank}{\operatorname{rank}}
\newcommand{\Span}{\operatorname{Span}}
\newcommand{\Ker}{\operatorname{Ker}}
\newcommand{\Imop}{\operatorname{Im}}
\newcommand{\diag}{\operatorname{diag}}
\newcommand{\Null}{\operatorname{Null}}

\newcommand{\nullity}{\operatorname{nullity}}

\theoremstyle{definition}
\newtheorem{definition}{Definition}[section]
\newtheorem{example}[definition]{Example}
\newtheorem{exercise}[definition]{Exercise}
\newtheorem{remark}[definition]{Remark}

\theoremstyle{plain}
\newtheorem{theorem}[definition]{Theorem}
\newtheorem{proposition}[definition]{Proposition}
\newtheorem{lemma}[definition]{Lemma}
\newtheorem{corollary}[definition]{Corollary}

\title{Eigenvectors and Diagonalization\\\large Full Lecture Notes}
\author{Cleaned-up LaTeX notes from lectures dated 2019/11/05, 2019/11/07, and 2019/11/08}
\date{}

\begin{document}
\maketitle
\tableofcontents

\section{Eigenvectors and Diagonalization I}

\subsection{The basic idea}

Let $V$ be a finite-dimensional vector space over a field $\F$, and let
\[
T:V\to V
\]
be a linear operator.

The main idea of diagonalization is this: even if $T$ looks complicated in one basis, it may look very simple in another basis.  In particular, suppose there is a basis
\[
\beta=\{u_1,\dots,u_n\}
\]
of $V$ such that
\[
[T]_{\beta}=\begin{pmatrix}
 d_1&0&\cdots&0\\
 0&d_2&\cdots&0\\
 \vdots&\vdots&\ddots&\vdots\\
 0&0&\cdots&d_n
\end{pmatrix}.
\]
Then
\[
T(u_i)=d_i u_i,\qquad 1\leq i\leq n.
\]
Thus, with respect to the basis $\beta$, the operator $T$ acts on each basis vector simply by scalar multiplication.

\subsection{A motivating example: reflection in $\R^2$}

Let $m\in\R$ and define $T:\R^2\to\R^2$ by
\[
T\begin{pmatrix}x\\y\end{pmatrix}
=\frac{1}{1+m^2}
\begin{pmatrix}
1-m^2&2m\\
2m&m^2-1
\end{pmatrix}
\begin{pmatrix}x\\y\end{pmatrix}.
\]
This is the reflection across the line spanned by $(1,m)^T$.

Let
\[
u_1=\begin{pmatrix}1\\m\end{pmatrix},
\qquad
u_2=\begin{pmatrix}m\\-1\end{pmatrix}.
\]
Then $\beta=\{u_1,u_2\}$ is a basis of $\R^2$, and direct computation gives
\[
T(u_1)=u_1,
\qquad
T(u_2)=-u_2.
\]
Therefore
\[
[T]_{\beta}=\begin{pmatrix}1&0\\0&-1\end{pmatrix}.
\]
So the same operator which may look complicated in the standard basis becomes extremely simple in the basis $\beta$.

Geometrically, the direction $\Span(u_1)$ is fixed by the reflection, and the perpendicular direction $\Span(u_2)$ is reversed.

\subsection{Eigenvectors and eigenvalues}

\begin{definition}[Eigenvector and eigenvalue]
Let $V$ be a finite-dimensional vector space over $\F$, and let $T:V\to V$ be a linear operator.

A nonzero vector $u\in V$ is called an \emph{eigenvector} of $T$ if there exists $\lambda\in\F$ such that
\[
T(u)=\lambda u.
\]
The scalar $\lambda$ is called the \emph{eigenvalue} corresponding to $u$.
\end{definition}

Equivalently,
\[
T(u)=\lambda u
\quad\Longleftrightarrow\quad
(T-\lambda\id_V)(u)=0,
\]
where $u\neq 0$.

\begin{remark}
If $u$ is an eigenvector of $T$ with eigenvalue $\lambda$, then the one-dimensional subspace
\[
\F u=\{cu:c\in\F\}
\]
is $T$-invariant, and the restriction
\[
T|_{\F u}:\F u\to\F u
\]
acts as scalar multiplication by $\lambda$:
\[
T(cu)=cT(u)=c\lambda u=\lambda(cu).
\]
\end{remark}

\begin{definition}[Diagonalizable linear operator]
A linear operator $T:V\to V$ is called \emph{diagonalizable} if there exists a basis $\beta$ of $V$ such that $[T]_{\beta}$ is a diagonal matrix.
\end{definition}

Equivalently,
\[
T\text{ is diagonalizable}
\quad\Longleftrightarrow\quad
\text{there exists a basis of }V\text{ consisting of eigenvectors of }T.
\]

Indeed, if
%\[
%\beta=\{u_1,\dots,u_n\}
%\]
is a basis of eigenvectors and
\[
T(u_i)=\lambda_i u_i,
\]
then
\[
[T]_{\beta}=\diag(\lambda_1,\dots,\lambda_n).
\]

\subsection{Matrix version of diagonalization}

Let
\[
A\in \Mat_{n\times n}(\F)
\]
and consider the linear operator
\[
L_A:\F^n\to\F^n,
\qquad
x\mapsto Ax.
\]

If $L_A$ is diagonalizable, then there exists a basis
\[
\beta=\{v_1,\dots,v_n\}
\]
of $\F^n$ consisting of eigenvectors of $A$.  Suppose
\[
Av_i=\lambda_i v_i.
\]
Let
\[
P=\begin{pmatrix} |&|&&|\\ v_1&v_2&\cdots&v_n\\ |&|&&|\end{pmatrix}.
\]
Then $P$ is invertible and
\[
AP=P\diag(\lambda_1,\dots,\lambda_n).
\]
Therefore
\[
P^{-1}AP=\diag(\lambda_1,\dots,\lambda_n).
\]

\begin{definition}[Diagonalizable matrix]
A matrix $A\in\Mat_{n\times n}(\F)$ is called \emph{diagonalizable over $\F$} if there exist an invertible matrix $P\in\Mat_{n\times n}(\F)$ and a diagonal matrix $D$ such that
\[
D=P^{-1}AP.
\]
Equivalently, $A$ is similar over $\F$ to a diagonal matrix.
\end{definition}

The columns of $P$ are eigenvectors of $A$, and the diagonal entries of $D$ are the corresponding eigenvalues.

\subsection{Example: rotation matrices and dependence on the field}

Let
\[
R_\theta=L_A,
\qquad
A=\begin{pmatrix}
\cos\theta&-\sin\theta\\
\sin\theta&\cos\theta
\end{pmatrix},
\qquad
0<\theta<2\pi,
\quad
\theta\neq\pi.
\]
This is rotation by angle $\theta$ in the plane.

\subsubsection*{Case 1: over $\R$}

We look for a nonzero vector $u\in\R^2$ and $\lambda\in\R$ such that
\[
Au=\lambda u.
\]
This is equivalent to
\[
(A-\lambda I)u=0,
\qquad u\neq 0.
\]
Thus $A-\lambda I$ must not be invertible, so
\[
\det(A-\lambda I)=0.
\]
Compute
\[
\det(A-\lambda I)
=\begin{vmatrix}
\cos\theta-\lambda&-\sin\theta\\
\sin\theta&\cos\theta-\lambda
\end{vmatrix}
=(\cos\theta-\lambda)^2+\sin^2\theta.
\]
Hence
\[
\lambda^2-2(\cos\theta)\lambda+1=0.
\]
The discriminant is
\[
\Delta=4\cos^2\theta-4=4(\cos^2\theta-1)=-4\sin^2\theta<0,
\]
because $\theta\neq 0,\pi,2\pi$.  Therefore there is no real eigenvalue.  Hence $R_\theta$ is not diagonalizable over $\R$.

\subsubsection*{Case 2: over $\C$}

Over $\C$, the equation
\[
\lambda^2-2(\cos\theta)\lambda+1=0
\]
has two distinct roots:
\[
\lambda_1=\cos\theta+i\sin\theta=e^{i\theta},
\qquad
\lambda_2=\cos\theta-i\sin\theta=e^{-i\theta}.
\]

For $\lambda_1=e^{i\theta}$,
\[
A-\lambda_1 I
=\begin{pmatrix}
-i\sin\theta&-\sin\theta\\
\sin\theta&-i\sin\theta
\end{pmatrix}.
\]
Solving
\[
(A-\lambda_1 I)
\begin{pmatrix}u_1\\u_2\end{pmatrix}=0
\]
gives
\[
-i u_1-u_2=0,
\qquad
u_2=-iu_1.
\]
So an eigenvector for $\lambda_1$ is
\[
\begin{pmatrix}1\\-i\end{pmatrix}.
\]

Similarly, for $\lambda_2=e^{-i\theta}$, an eigenvector is
\[
\begin{pmatrix}1\\i\end{pmatrix}.
\]
Thus
\[
\beta=
\left\{
\begin{pmatrix}1\\-i\end{pmatrix},
\begin{pmatrix}1\\i\end{pmatrix}
\right\}
\]
is a basis of $\C^2$, and
\[
[R_\theta]_{\beta}
=\begin{pmatrix}
e^{i\theta}&0\\
0&e^{-i\theta}
\end{pmatrix}.
\]
Therefore the rotation matrix is diagonalizable over $\C$, but not over $\R$.

\subsection{Characteristic polynomial of a matrix}

\begin{definition}[Characteristic polynomial]
Let $A\in\Mat_{n\times n}(\F)$.  The characteristic polynomial of $A$ is
\[
\varphi_A(t)=\det(A-tI).
\]
\end{definition}

A scalar $\lambda\in\F$ is an eigenvalue of $A$ if and only if
\[
\varphi_A(\lambda)=0.
\]

\begin{exercise}[Coefficients of the characteristic polynomial]
Show that
\[
\varphi_A(t)=(-1)^n\left(t^n+a_1t^{n-1}+\cdots+a_{n-1}t+a_n\right),
\]
where
\[
a_1=-\tr(A),
\qquad
(-1)^n a_n=\det(A).
\]
\end{exercise}

\begin{proof}[Solution sketch]
The constant term of $\varphi_A(t)=\det(A-tI)$ is obtained by setting $t=0$, so it is $\det(A)$.  Thus
\[
(-1)^n a_n=\det(A).
\]

The highest degree term comes from multiplying the diagonal terms $-t$, giving $(-t)^n=(-1)^n t^n$.  The coefficient of $t^{n-1}$ comes from choosing exactly one diagonal entry $a_{ii}$ and choosing $-t$ from the remaining diagonal positions.  Hence that coefficient is
\[
(-1)^{n-1}\sum_{i=1}^n a_{ii}=(-1)^{n-1}\tr(A).
\]
Comparing with
\[
(-1)^n a_1 t^{n-1},
\]
we get
\[
a_1=-\tr(A).
\]
\end{proof}

\section{Eigenvectors and Diagonalization II}

\subsection{Determinant and trace of a linear operator}

Let $T:V\to V$ be a linear operator on a finite-dimensional vector space $V$ over $\F$.

\begin{definition}[Determinant and trace of a linear operator]
Choose any basis $\beta$ of $V$.  Define
\[
\det(T):=\det([T]_{\beta}),
\qquad
\tr(T):=\tr([T]_{\beta}).
\]
\end{definition}

We must check that these definitions do not depend on the choice of basis.

\begin{lemma}[Similarity preserves determinant and trace]
Let $A,B\in\Mat_{n\times n}(\F)$.  Suppose there exists an invertible matrix $P$ such that
\[
A=P^{-1}BP.
\]
Then
\[
\det(A)=\det(B),
\qquad
\tr(A)=\tr(B).
\]
\end{lemma}

\begin{proof}
For the determinant,
\[
\det(A)=\det(P^{-1}BP)
=\det(P^{-1})\det(B)\det(P)
=\det(B).
\]
For the trace, use the identity
\[
\tr(XY)=\tr(YX)
\]
for square matrices $X,Y$.  Then
\[
\tr(A)=\tr(P^{-1}BP)
=\tr(BPP^{-1})
=\tr(B).
\]
\end{proof}

If $\beta$ and $\gamma$ are two bases of $V$, then $[T]_{\beta}$ and $[T]_{\gamma}$ are similar matrices.  Therefore determinant and trace are well-defined invariants of the operator $T$.

\subsection{Eigenvalues and the determinant condition}

Recall that $u\in V$ is an eigenvector of $T$ with eigenvalue $\lambda$ if and only if
\[
u\neq 0,
\qquad
(T-\lambda\id_V)(u)=0.
\]
Thus
\[
\lambda\text{ is an eigenvalue}
\quad\Longleftrightarrow\quad
\Ker(T-\lambda\id_V)\neq\{0\}.
\]

\begin{proposition}
Let $T:V\to V$ be a linear operator on a finite-dimensional vector space.  Then
\[
\lambda\in\F\text{ is an eigenvalue of }T
\quad\Longleftrightarrow\quad
\det(T-\lambda\id_V)=0.
\]
\end{proposition}

\begin{proof}
Choose a basis $\beta$ of $V$.  Then
\[
[T-\lambda\id_V]_{\beta}=[T]_{\beta}-\lambda I.
\]
Now
\[
\Ker(T-\lambda\id_V)\neq\{0\}
\]
if and only if the matrix $[T]_{\beta}-\lambda I$ has a nontrivial kernel.  Since $V$ is finite-dimensional, this happens if and only if the matrix is not invertible, which happens if and only if its determinant is zero.  Hence
\[
\lambda\text{ is an eigenvalue}
\quad\Longleftrightarrow\quad
\det([T]_{\beta}-\lambda I)=0.
\]
But this determinant is exactly $\det(T-\lambda\id_V)$.
\end{proof}

\begin{definition}[Characteristic polynomial of a linear operator]
The characteristic polynomial of $T$ is
\[
\varphi_T(t)=\det(T-t\id_V).
\]
\end{definition}

Therefore,
\[
\lambda\text{ is an eigenvalue of }T
\quad\Longleftrightarrow\quad
\varphi_T(\lambda)=0.
\]

\subsection{Linear independence of eigenvectors with distinct eigenvalues}

\begin{lemma}
Let $u_1,\dots,u_k$ be eigenvectors of $T$ with eigenvalues
\[
\lambda_1,\dots,\lambda_k.
\]
Assume the eigenvalues are mutually distinct:
\[
\lambda_i\neq\lambda_j
\quad\text{whenever}\quad i\neq j.
\]
Then
\[
\{u_1,\dots,u_k\}
\]
is linearly independent.
\end{lemma}

\begin{proof}
Suppose not.  Let $h$ be the smallest positive integer such that
\[
\{u_1,\dots,u_h\}
\]
is linearly dependent.  Then $h\geq 2$, and $\{u_1,\dots,u_{h-1}\}$ is linearly independent.

There exist scalars $a_1,\dots,a_h$, not all zero, such that
\[
a_1u_1+a_2u_2+\cdots+a_hu_h=0. \tag{1}
\]
By minimality of $h$, we must have $a_h\neq 0$.

Apply $T$ to (1).  Since $T(u_i)=\lambda_i u_i$, we get
\[
a_1\lambda_1u_1+a_2\lambda_2u_2+\cdots+a_h\lambda_hu_h=0. \tag{2}
\]
Multiply (1) by $\lambda_h$:
\[
a_1\lambda_hu_1+a_2\lambda_hu_2+\cdots+a_h\lambda_hu_h=0. \tag{3}
\]
Subtract (3) from (2):
\[
a_1(\lambda_1-\lambda_h)u_1+\cdots+a_{h-1}(\lambda_{h-1}-\lambda_h)u_{h-1}=0.
\]
Since $u_1,\dots,u_{h-1}$ are linearly independent, each coefficient is zero:
\[
a_i(\lambda_i-\lambda_h)=0,
\qquad 1\leq i\leq h-1.
\]
Because the eigenvalues are distinct, $\lambda_i-\lambda_h\neq 0$, so
\[
a_i=0,
\qquad 1\leq i\leq h-1.
\]
Then (1) becomes
\[
a_hu_h=0.
\]
Since $u_h\neq 0$, we get $a_h=0$, contradiction.

Therefore the eigenvectors are linearly independent.
\end{proof}

\begin{corollary}
Let $\dim V=n$.  If $T:V\to V$ has $n$ distinct eigenvalues in $\F$, then $T$ is diagonalizable over $\F$.
\end{corollary}

\begin{proof}
Choose one eigenvector for each eigenvalue.  By the lemma, these $n$ eigenvectors are linearly independent.  Since $\dim V=n$, they form a basis of $V$.  Hence $V$ has a basis consisting of eigenvectors, so $T$ is diagonalizable.
\end{proof}

\begin{remark}
The converse is false.  A linear operator may be diagonalizable even if it has fewer than $n$ distinct eigenvalues.  The identity map on an $n$-dimensional vector space has only one eigenvalue, $1$, but it is diagonalizable.
\end{remark}

\section{Diagonalization III}

\subsection{Eigenspaces}

\begin{definition}[Eigenspace]
Let $\lambda$ be an eigenvalue of $T:V\to V$.  Define
\[
E_{\lambda}=\{u\in V:T(u)=\lambda u\}.
\]
This is called the \emph{eigenspace} corresponding to $\lambda$.
\end{definition}

Since
\[
T(u)=\lambda u
\quad\Longleftrightarrow\quad
(T-\lambda\id_V)(u)=0,
\]
we have
\[
E_{\lambda}=\Ker(T-\lambda\id_V).
\]
Therefore $E_{\lambda}$ is a subspace of $V$.

\begin{exercise}
Show directly that $E_{\lambda}\leq V$.
\end{exercise}

\begin{proof}[Solution]
Let $u,v\in E_{\lambda}$ and $c\in\F$.  Then
\[
T(u+v)=T(u)+T(v)=\lambda u+\lambda v=\lambda(u+v),
\]
so $u+v\in E_{\lambda}$.  Also
\[
T(cu)=cT(u)=c\lambda u=\lambda(cu),
\]
so $cu\in E_{\lambda}$.  Thus $E_{\lambda}$ is a subspace.
\end{proof}

\subsection{Rank-nullity and eigenspace dimension}

For a linear map $T:V\to W$, the first isomorphism theorem gives
\[
V/\Ker(T)\cong \Imop(T).
\]
Taking dimensions gives
\[
\dim V-\dim\Ker(T)=\dim\Imop(T).
\]
Therefore
\[
\dim V=\nullity(T)+\rank(T).
\]
This is the rank-nullity theorem.

Applying this to $T-\lambda\id_V$, we get
\[
\dim E_{\lambda}
=\dim\Ker(T-\lambda\id_V)
=\dim V-\rank(T-\lambda\id_V).
\]

The number
\[
\dim E_{\lambda}
\]
is called the \emph{geometric multiplicity} of the eigenvalue $\lambda$.

\subsection{Example: a diagonalizable $2\times2$ matrix}

Let
\[
A=\begin{pmatrix}1&1\\1&1\end{pmatrix}.
\]
Then
\[
\varphi_A(t)=\det(A-tI)
=\begin{vmatrix}1-t&1\\1&1-t\end{vmatrix}
=(1-t)^2-1
=t^2-2t
=t(t-2).
\]
So the eigenvalues are
\[
\lambda=0,
\qquad
\lambda=2.
\]

For $\lambda=0$,
\[
E_0=\Ker(A)=\Span\left\{\begin{pmatrix}1\\-1\end{pmatrix}\right\}.
\]
For $\lambda=2$,
\[
E_2=\Ker(A-2I)=\Span\left\{\begin{pmatrix}1\\1\end{pmatrix}\right\}.
\]
Thus $A$ has two linearly independent eigenvectors in $\R^2$, so it is diagonalizable.

\subsection{Example: diagonalizable with repeated eigenvalue}

Let
\[
A=\begin{pmatrix}
4&0&1\\
2&3&2\\
1&0&4
\end{pmatrix}.
\]
Then
\[
\varphi_A(t)=\det(A-tI)
=\begin{vmatrix}
4-t&0&1\\
2&3-t&2\\
1&0&4-t
\end{vmatrix}.
\]
Expanding along the second column,
\[
\varphi_A(t)
=(3-t)\begin{vmatrix}4-t&1\\1&4-t\end{vmatrix}
=(3-t)((4-t)^2-1).
\]
Thus
\[
\varphi_A(t)=(3-t)(t^2-8t+15)
=(3-t)(t-3)(t-5)
=-(t-3)^2(t-5).
\]
The eigenvalues are
\[
3,
\qquad
5.
\]
There are only two distinct eigenvalues, but $A$ may still be diagonalizable.

For $\lambda=5$,
\[
A-5I=
\begin{pmatrix}
-1&0&1\\
2&-2&2\\
1&0&-1
\end{pmatrix}.
\]
The rank is $2$, so
\[
\dim E_5=3-2=1.
\]
Solving gives
\[
E_5=\Span\left\{
\begin{pmatrix}1\\2\\1\end{pmatrix}
\right\}.
\]

For $\lambda=3$,
\[
A-3I=
\begin{pmatrix}
1&0&1\\
2&0&2\\
1&0&1
\end{pmatrix}.
\]
The rank is $1$, so
\[
\dim E_3=3-1=2.
\]
Solving gives
\[
E_3=\Span\left\{
\begin{pmatrix}1\\0\\-1\end{pmatrix},
\begin{pmatrix}0\\1\\0\end{pmatrix}
\right\}.
\]

Hence
\[
\dim E_3+\dim E_5=2+1=3=\dim\R^3.
\]
Therefore $A$ is diagonalizable.  One eigenbasis is
\[
\beta=
\left\{
\begin{pmatrix}1\\0\\-1\end{pmatrix},
\begin{pmatrix}0\\1\\0\end{pmatrix},
\begin{pmatrix}1\\2\\1\end{pmatrix}
\right\},
\]
and with respect to this basis,
\[
[L_A]_{\beta}=
\begin{pmatrix}
3&0&0\\
0&3&0\\
0&0&5
\end{pmatrix}.
\]

\subsection{Sums and direct sums of subspaces}

\begin{definition}[Sum of subspaces]
Let $U,W\leq V$.  Define
\[
U+W=\{u+w:u\in U,\ w\in W\}.
\]
This is called the sum of $U$ and $W$.
\end{definition}

\begin{exercise}
Show that
\[
U+W\leq V
\]
and
\[
U\cap W\leq U+W.
\]
\end{exercise}

\begin{definition}[Direct sum]
Let $U,W\leq V$.  We say
\[
V=U\oplus W
\]
if
\[
V=U+W
\]
and
\[
U\cap W=\{0\}.
\]
\end{definition}

Equivalently,
\[
V=U\oplus W
\]
means that every vector $v\in V$ can be written uniquely as
\[
v=u+w,
\qquad
u\in U,
\quad
w\in W.
\]

More generally, for subspaces $W_1,\dots,W_m\leq V$, we write
\[
V=W_1\oplus W_2\oplus\cdots\oplus W_m
\]
if every vector $v\in V$ can be written uniquely as
\[
v=w_1+w_2+\cdots+w_m,
\qquad
w_i\in W_i.
\]

Equivalently,
\[
V=W_1+\cdots+W_m
\]
and
\[
W_j\cap\left(W_1+\cdots+W_{j-1}+W_{j+1}+\cdots+W_m\right)=\{0\}
\]
for each $j$.

\begin{exercise}
Let $u_1,\dots,u_n$ be a basis for $V$.  Let
\[
U=\Span\{u_1,\dots,u_k\},
\qquad
W=\Span\{u_{k+1},\dots,u_n\}.
\]
Show that
\[
V=U\oplus W.
\]
\end{exercise}

\subsection{Diagonalization criterion using eigenspaces}

Let $T:V\to V$ be a linear operator on an $n$-dimensional vector space.  Suppose the distinct eigenvalues of $T$ are
\[
\lambda_1,\dots,\lambda_k,
\qquad k\leq n.
\]
For each $i$, let
\[
E_{\lambda_i}=\Ker(T-\lambda_i\id_V).
\]

\begin{theorem}[Diagonalization theorem]
The linear operator $T$ is diagonalizable if and only if
\[
V=E_{\lambda_1}\oplus E_{\lambda_2}\oplus\cdots\oplus E_{\lambda_k}.
\]
\end{theorem}

\begin{proof}
First suppose $T$ is diagonalizable.  Then there exists a basis $\beta$ of $V$ consisting of eigenvectors.  Group this basis according to eigenvalue:
\[
\beta=
\{u_{11},\dots,u_{1m_1},
 u_{21},\dots,u_{2m_2},
 \dots,
 u_{k1},\dots,u_{km_k}\},
\]
where each $u_{ij}$ has eigenvalue $\lambda_i$.

Let
\[
W_i=\Span\{u_{i1},\dots,u_{im_i}\}.
\]
Since $\beta$ is a basis,
\[
V=W_1\oplus\cdots\oplus W_k.
\]
Clearly
\[
W_i\subseteq E_{\lambda_i}.
\]
We claim that actually $W_i=E_{\lambda_i}$.

Let $v\in E_{\lambda_i}$.  Since $V=W_1\oplus\cdots\oplus W_k$, write uniquely
\[
v=w_1+\cdots+w_k,
\qquad
w_j\in W_j.
\]
Then
\[
T(v)=\lambda_i v.
\]
But also, since $T$ acts as multiplication by $\lambda_j$ on $W_j$,
\[
T(v)=\lambda_1w_1+\cdots+\lambda_kw_k.
\]
Thus
\[
0=T(v)-\lambda_i v
=\sum_{j=1}^k(\lambda_j-\lambda_i)w_j.
\]
By uniqueness in the direct sum, each component must be zero.  Hence for $j\neq i$,
\[
(\lambda_j-\lambda_i)w_j=0.
\]
Since $\lambda_j\neq\lambda_i$, we get $w_j=0$ for all $j\neq i$.  Therefore $v=w_i\in W_i$.  Thus
\[
E_{\lambda_i}=W_i.
\]
Hence
\[
V=E_{\lambda_1}\oplus\cdots\oplus E_{\lambda_k}.
\]

Conversely, suppose
\[
V=E_{\lambda_1}\oplus\cdots\oplus E_{\lambda_k}.
\]
For each $i$, choose a basis $\beta_i$ of $E_{\lambda_i}$.  Then
\[
\beta=\beta_1\cup\cdots\cup\beta_k
\]
is a basis of $V$ because the sum is direct.  Every vector in $\beta_i$ is an eigenvector with eigenvalue $\lambda_i$.  Therefore $\beta$ is a basis of $V$ consisting of eigenvectors.  Hence $T$ is diagonalizable.
\end{proof}

\begin{corollary}[Dimension criterion]
Let $T:V\to V$ have distinct eigenvalues $\lambda_1,\dots,\lambda_k$ in $\F$.  Then $T$ is diagonalizable over $\F$ if and only if
\[
\dim E_{\lambda_1}+\cdots+\dim E_{\lambda_k}=\dim V.
\]
\end{corollary}

\begin{proof}
Eigenspaces corresponding to distinct eigenvalues have zero intersection in the direct-sum sense, by the linear-independence lemma.  Therefore
\[
E_{\lambda_1}+\cdots+E_{\lambda_k}
\]
is automatically a direct sum.  Thus it equals $V$ if and only if its dimension equals $\dim V$, namely if and only if
\[
\sum_{i=1}^k\dim E_{\lambda_i}=\dim V.
\]
\end{proof}

\begin{remark}
If
\[
V=E_{\lambda_1}\oplus\cdots\oplus E_{\lambda_k}
\]
and $m_i=\dim E_{\lambda_i}$, then with respect to a basis obtained by joining bases of the eigenspaces,
\[
[T]_{\beta}
=\begin{pmatrix}
\lambda_1 I_{m_1}&&&0\\
&\lambda_2 I_{m_2}&&\\
&&\ddots&\\
0&&&\lambda_k I_{m_k}
\end{pmatrix}.
\]
This is a diagonal matrix whose diagonal entries are the eigenvalues repeated according to the dimensions of their eigenspaces.
\end{remark}

\section{Summary}

\begin{itemize}
\item An eigenvector is a nonzero vector $u$ satisfying $T(u)=\lambda u$.
\item A linear operator is diagonalizable if and only if it has a basis consisting of eigenvectors.
\item For a matrix $A$, diagonalization means finding $P$ invertible and $D$ diagonal such that
\[
D=P^{-1}AP.
\]
\item The eigenvalues of $A$ are the roots of
\[
\varphi_A(t)=\det(A-tI).
\]
\item Eigenvectors corresponding to distinct eigenvalues are linearly independent.
\item If $\dim V=n$ and $T$ has $n$ distinct eigenvalues, then $T$ is diagonalizable.
\item The eigenspace corresponding to $\lambda$ is
\[
E_\lambda=\Ker(T-\lambda\id_V).
\]
\item The operator $T$ is diagonalizable if and only if
\[
V=\bigoplus_{\lambda} E_\lambda,
\]
where the sum runs over the distinct eigenvalues of $T$.
\item Equivalently,
\[
T\text{ is diagonalizable}
\quad\Longleftrightarrow\quad
\sum_{\lambda}\dim E_\lambda=\dim V,
\]
provided all eigenvalues needed for diagonalization lie in the field $\F$.
\item Diagonalizability depends on the field.  A real rotation matrix with angle not $0$ or $\pi$ is not diagonalizable over $\R$, but it is diagonalizable over $\C$.
\end{itemize}

\end{document}
